\chapter{One Cosine Primitive, Three Proofs}\label{ch:three-cosine-proofs}

\begin{center}
\href{cosine-primitive-graph.html}{\textbf{Open the focused three-proof dependency graph}}
\end{center}

The theorem occurs once. Three alternative derivations end at that node:
finite inequalities, native concave FTC, and Mathlib integration with
representation bridges. The graph selects mathematical milestones, not every
helper declaration. A shared rational root branches into native interval
computations and Mathlib real numbers. Three labeled proof lanes keep the
arguments apart; genuine common constructions are shown once. All three
complete Lean theorem types are checked to be definitionally identical.

\begin{definition}[The shared rational numbers]\label{def:c3-rationals}
\lean{Rat,Rat.num,Rat.den,Rat.normalize,Rat.add,Rat.mul}\leanok
Both foundations start with Lean's exact rational type \code{Rat}, written
$\mathbb{Q}$. It supplies normalized fractions and exact rational arithmetic.
Our finite interval algorithms do not need Mathlib's completed real field.
\end{definition}

\begin{definition}[Mathlib's real-number construction]\label{def:c3-mreal}
\uses{def:c3-rationals}
\lean{Real,Real.ofCauchy,CauSeq.Completion.Cauchy,Real.cauchy,Real.mk,Real.mk_eq}\mathlibok
Mathlib's real type is the Cauchy completion of these same rationals.
\code{Real.mk} maps a rational Cauchy sequence to its real value, and
\code{Real.mk\_eq} identifies equality with equivalence of representatives.
This is distinct from the comparison interpretation below, which relates
our nested intervals to an already constructed Mathlib real.
\end{definition}

\begin{definition}[Rational interval computations]\label{def:c3-intervals}
\uses{def:c3-rationals}
\lean{ComputableAnalysis.RealRaw.Valid,ComputableAnalysis.RealRaw.Equiv}\leanok
A valid raw real is a program producing ordered nested rational intervals
whose widths tend to zero. Equivalence means stagewise overlap. Arithmetic
and finite intersections are rational computations; equivalence becomes
transitive for valid representatives.
\end{definition}

\begin{definition}[Arctangent before trigonometry]\label{def:c3-arctan}
\uses{def:c3-intervals}
\lean{ComputableAnalysis.CosinePrimitive.A,ComputableAnalysis.ArctanGeometry.arctanIntegralRectangleRaw_equiv_arctanGeom}\leanok
For $0\le u\le1$, use the geometric arctangent program $A(u)$, with its
proved equivalence to the rational rectangle computation
\[
 A(u)\simeq\int_0^u\frac{dv}{1+v^2}.
\]
The integrand is rational; no sine or cosine is used to construct $A$.
\end{definition}

\begin{definition}[The arctangent presentation of pi]\label{def:c3-pi}
\uses{def:c3-arctan}
\lean{ComputableAnalysis.CosinePrimitive.pi,ComputableAnalysis.CosinePrimitive.pi_compute}\leanok
Define the program $\pi:=4A(1)$. It agrees stage by stage with the existing
circle-area presentation. The reciprocal computation is certified positive;
other valid pi presentations can be transported through this equivalence.
\end{definition}

\begin{lemma}[The closed inverse provider]\label{lem:c3-inverse}
\uses{def:c3-arctan}
\lean{ComputableAnalysis.ClosedArctanInverse.provider}\leanok
An explicit finite inverse search on the unit-slope chart computes, for
rational $0\le x\le1/2$, a parameter $u(x)$ in $[0,1]$ with
$A(u(x))\simeq2xA(1)$. No caller-supplied inverse provider remains.
\end{lemma}
\begin{proof}\leanok
Rational bisection of a finite clock, widening, and prefix intersection give
nested source boxes of width at most $56/2^n$. A finite forward-error estimate
identifies their arctangent with the target. The grouped Lean declarations
show the evaluator, range, convergence, and forward-equation certificates.
\end{proof}

\begin{definition}[Our sine and cosine]\label{def:c3-trig}
\uses{lem:c3-inverse,def:c3-pi}
\lean{ComputableAnalysis.CosinePrimitive.S,ComputableAnalysis.CosinePrimitive.C}\leanok
Rational circle coordinates of the inverse parameter define
\[
 S(x)=\frac{2u(x)}{1+u(x)^2},\qquad C(x)=\frac{1-u(x)^2}{1+u(x)^2}.
\]
These formulas denote interval algorithms. The input is a rational normalized
angle. Their agreement with conventional sine and cosine is a separate
representation theorem, not their definition.
\end{definition}

\begin{definition}[The fixed cosine integral program]\label{def:c3-integral}
\uses{def:c3-trig,def:c3-intervals}
\lean{ComputableAnalysis.CosinePrimitive.integral}\leanok
Let $R_{k,n}$ be the cosine left rectangle sum on $k+1$ equal cells, with
samples evaluated at stage $n$. For rational $0\le t\le1/2$, define
\[
 I(t)_n=\bigcap_{k=0}^{n}
 \operatorname{expand}\left(R_{k,n},\frac{4000t^2}{k+1}\right).
\]
The program reads cosine samples, not sine endpoints. Retaining old meshes
allows each fixed finite sample set to converge without an assumed uniform
rate for newly introduced inputs. The error constant is conservative.
\end{definition}

\begin{lemma}[Direct finite inequalities]\label{lem:c3-direct}
\uses{def:c3-trig,def:c3-arctan}
\lean{ComputableAnalysis.GeometricSineDirectBounds.positive_increment_error,ComputableAnalysis.CosineFTC.fixedMesh_overlaps_endpoint}\leanok
For fixed positive rational $h$, sufficiently late rational selections satisfy
$|s_1-s_0-hpc|\le4000h^2$. Finite telescoping and a reciprocal-pi corner
compare each widened cosine sum with sine endpoints.
\end{lemma}
\begin{proof}\leanok
Use rational circle and clock residual identities, not a derivative or FTC.
Choose a common stage for finitely many cells; telescope rational selections
and transport to arbitrary stages by nesting. The value $S(0)\simeq0$
removes the lower endpoint in the basepoint statement.
\end{proof}

\begin{lemma}[Concavity and supporting secants]\label{lem:c3-concavity}
\uses{def:c3-trig,def:c3-arctan}
\lean{ComputableAnalysis.GeometricSineConcavity.primitive_concave,ComputableAnalysis.GeometricSineConcavity.primitiveDerivativeData}\leanok
For $P=S/\pi$ and rational $x<y$ in the chart,
\[
 C(y)\preceq\frac{P(y)-P(x)}{y-x}\preceq C(x).
\]
The geometric proof supplies exact concavity and a Lipschitz constant $32$
for the supporting slope computation. Neither derivative nor integral
identities are assumed to prove curvature.
\end{lemma}
\begin{proof}\leanok
Compare rational circle chords with tangent slopes, then use the arctangent
clock inequalities. Remove arbitrary rational slack by validity and nesting.
The same data proves a valid sine-only secant derivative at interior points;
that companion theorem is not itself a premise of the FTC.
\end{proof}

\begin{theorem}[Native concave FTC]\label{thm:c3-ftc}
\uses{def:c3-intervals}
\lean{ComputableAnalysis.ConcaveFTC.integral_equiv_endpoint}\leanok
A concave primitive with valid supporting secants and a quantitative slope
bound has its derivative integral identified with its endpoint difference.
The certificate contains no assumed integral equivalence.
\end{theorem}
\begin{proof}\leanok
When the relevant source boxes have width at most $\tau$, the local estimate
is $|F_n^-(y)-F_n^-(x)-hD_n^-(x)|\le Kh^2+(1+h)\tau$.
For a fixed finite mesh, select a common stage, telescope, remove arbitrary
positive slack, and intersect the resulting enclosures.
\end{proof}

\begin{definition}[Interpretation in Mathlib reals]\label{def:c3-real}
\uses{def:c3-intervals,def:c3-mreal}
\lean{MathlibComparison.Represents,MathlibComparison.equiv_of_represents}\leanok
The optional comparison sets $\operatorname{Represents}(X,r)$ to mean that
every rational output box of $X$ contains the Mathlib real $r$. Two programs
representing the same value overlap at every stage. The current general
interpretation uses a supremum in the comparison layer; the native algorithms
do not acquire this dependency.
\end{definition}

\begin{definition}[Mathlib exponential and trigonometry]\label{def:c3-mexp}
\uses{def:c3-mreal}
\lean{Complex.exp,Real.sin,Real.cos}\mathlibok
Mathlib constructs its complex exponential independently of our evaluator,
then defines sine and cosine from exponentials. These are target functions
of the representation bridge, not replacements for our programs.
\end{definition}

\begin{definition}[Mathlib pi]\label{def:c3-mpi}
\uses{def:c3-mexp}
\lean{Real.pi}\mathlibok
Mathlib defines pi as twice a chosen cosine zero in $[1,2]$.
The identity $\arctan(1)=\pi/4$ bridges this to our $4A(1)$ without
unfolding that choice in every subsequent comparison.
\end{definition}

\begin{lemma}[Pi, sine, and cosine value bridges]\label{lem:c3-values}
\uses{def:c3-real,def:c3-trig,def:c3-pi,def:c3-mexp,def:c3-mpi}
\lean{MathlibComparison.piCircleArea_represents,MathlibComparison.sine_represents,MathlibComparison.cosine_represents}\leanok
Our pi represents Mathlib pi. For every rational $x$ in the chart, $S(x)$
represents $\operatorname{Real.sin}(\operatorname{Real.pi}\,x)$, while $C(x)$
represents the corresponding Mathlib cosine.
\end{lemma}
\begin{proof}\leanok
Compare arctangent rectangles with Mathlib's rational-kernel integral.
Interpret the inverse through convergent rational approximations, then apply
double-arctangent coordinate identities. Neither native endpoint proof is used.
\end{proof}

\begin{lemma}[The actual quadrature represents the Mathlib integral]\label{lem:c3-quadrature}
\uses{lem:c3-values,def:c3-integral,def:c3-real}
\lean{MathlibComparison.cosine_integral_represents}\leanok
Our $I(t)$ represents Mathlib's interval integral of
$g(x)=\operatorname{Real.cos}(\operatorname{Real.pi}\,x)$ from zero to $t$.
\end{lemma}
\begin{proof}\leanok
Each finite rational-box sum contains the corresponding real sampled sum.
The estimate $|g(x)-g(y)|\le4|x-y|$ bounds the rectangle error on $m$ cells
by $4t^2/m$, below our radius $4000t^2/m$. Integrate bounds on each cell
and add over the partition. No primitive theorem is used at this step.
Every expanded enclosure, hence its finite intersection, contains the same
Mathlib integral. This requires more than agreement at rational inputs alone.
\end{proof}

\begin{lemma}[Mathlib cosine primitive]\label{lem:c3-mftc}
\uses{def:c3-mexp,def:c3-mpi}
\lean{MathlibComparison.mathlib_cosine_primitive}\leanok
For every Mathlib real $t$,
\[
 \int_0^t\operatorname{Real.cos}(\operatorname{Real.pi}\,x)\,dx
 =\frac{\operatorname{Real.sin}(\operatorname{Real.pi}\,t)}{\operatorname{Real.pi}}.
\]
\end{lemma}
\begin{proof}\leanok
Apply Mathlib's cosine integral and linear change of variable. This becomes
a proof of our computational statement only with both bridges above.
\end{proof}

\begin{theorem}[One basepoint statement]\label{thm:c3-primitive}
\uses{def:c3-integral,def:c3-trig,def:c3-pi}
\lean{ComputableAnalysis.CosinePrimitive.Statement}\leanok
For rational $0\le t\le1/2$,
\[
 \boxed{\int_0^t C(x)\,dx\simeq\frac{S(t)}{\pi}.}
\]
The left side is the fixed cosine-sum program, not primitive evaluation.
The right side independently multiplies sine by reciprocal pi.
\end{theorem}
\begin{proof}[Three alternative proofs]
\lean{ComputableAnalysis.CosinePrimitive.viaInequalities,ComputableAnalysis.CosinePrimitive.viaFTC,ComputableAnalysis.CosinePrimitive.viaMathlib}\leanok
\uses{lem:c3-direct,lem:c3-concavity,thm:c3-ftc,lem:c3-quadrature,lem:c3-values,lem:c3-mftc}
\emph{Direct inequalities:} use the finite residual and telescope, then
$S(0)\simeq0$. \emph{Native FTC:} instantiate the concave theorem with
$P=S/\pi$ and $D=C$, and remove the lower endpoint.
\emph{Mathlib:} combine the independent quadrature correspondence, primitive
identity, and value bridges. Both sides represent the same real, so their
rational boxes overlap.

These are alternatives, not joint prerequisites. Stock blueprint metadata
stores one proof attachment per statement; this block records their union,
while the focused graph retains separate routes. The complete theorem types
and proof separation are audited. The native routes have no Mathlib
references; the third has an independent quadrature validity proof.
\end{proof}

\begin{lemma}[Companion: a native computable complex exponential]\label{lem:c3-native-exp}
\uses{def:c3-intervals}
\lean{ComputableAnalysis.RotationSeries.rotationExpRaw_valid}\leanok
The project has a factorial-series complex-box computation for $e^{iT}$ at
rational $T$, with proved validity and explicit tail bounds. It is native
and computable, not Mathlib's exponential.
\end{lemma}
\begin{proof}\leanok
Finite rational complex sums and factorial tail enclosures supply the
certificate. This is not a premise of the three integral proofs. Extension
to arbitrary represented complex inputs and identification with geometric
trigonometry are separate tasks, not asserted here. The graph displays this
companion only on request.
\end{proof}

\section{Display and quantitative comparison}
The final application, native prerequisites, comparison bridges, and Mathlib
prerequisites are measured separately. Stored types and proof bodies, not
import lists, determine dependencies. Explicit shared constructions remain
visible above the separate proof lanes; visual separation does not pretend
that the mathematical foundations are disjoint.

Each broad node opens grouped exact Lean declarations. Lexical syntax colors
highlight keywords, declaration names, types, literals, and operators without
changing a single character of the displayed or copied statement. Node and
card backgrounds separately indicate dependence on Mathlib's real type.
The one conclusion has three proof terms; only the third uses Mathlib reals.
